This paper studies zero-error design patterns in Python: patterns that remove ambiguity, prevent stale derived state, and make structural facts explicit. The OpenHCS-derived packages serve as the concrete evidence base, but the patterns themselves are general and reusable rather than OpenHCS-specific.

The result is a catalog of 13 recurring instances, organized around five headline mechanisms: provenance-preserving dual-axis resolution, branching state DAGs, zero-boilerplate registry and discovery, collision-safe executable serialization, and transport/backend separation. The broader theory from Papers 1--4 explains why these patterns matter; this paper focuses on the pattern level itself and on the proof obligations they satisfy in real code.
